% GENERATED FILE. Edit canonical lessons, then run npm run generate:books.
% canonical-source-hash: fnv1a64:6a23aa47d64e9236
% canonical-lessons: TE-C313-suddamukka, TE-C313-palaka, TE-C313-balapam, TE-C313-sira, TE-C313-hajaru

\chapter{Piece of chalk, Slate, Slate pencil, Ink, Attendance}
\label{ch:te-piece-of-chalk-slate-slate-pencil-ink-attendance}
\hlchaptermodality{313}

\begin{chapteropening}
\textbf{By the end of this chapter:} \emph{I can say a piece of chalk, a slate, a slate pencil, ink and attendance.}

Complete the last lesson of chapter 313: I can say a piece of chalk, a slate, a slate pencil, ink and attendance.
\end{chapteropening}

\section[\texorpdfstring{\te{సుద్దముక్క} (suddamukka)}{సుద్దముక్క (suddamukka)}]{\te{సుద్దముక్క} (suddamukka) — a piece of chalk}

\label{lesson:TE-C313-suddamukka}



\begin{quote}
Before the new one: say the Telugu for a coach, then the Telugu for an engineer.
\end{quote}

\subsection*{You'll want to know: \te{సుద్దముక్క}}
\textbf{\te{సుద్దముక్క}} — \emph{suddamukka} — \textquotedblleft{}a piece of chalk\textquotedblright{}.

\begin{grammarlens}[title={What you've built}]
One more word for this chapter.
\end{grammarlens}

\subsection*{Your turn}
\begin{itemize}
\raggedright
  \item \emph{Say it:} \emph{suddamukka}
  \item \emph{Say it:} \emph{suddamukka}, once more
  \item \emph{Say it:} say \emph{suddamukka}
  \item \emph{Recall:} say the Telugu for a coach, then the Telugu for an engineer, then say \emph{suddamukka} again
\end{itemize}

\begin{tcolorbox}[breakable,colback=teal!4,colframe=teal!35!black,title={Before you move on}]
What does \te{సుద్దముక్క} mean? (\textquotedblleft{}A piece of chalk\textquotedblright{}.) Say it once more.
\end{tcolorbox}
\section[\texorpdfstring{\te{పలక} (palaka)}{పలక (palaka)}]{\te{పలక} (palaka) — a slate}

\label{lesson:TE-C313-palaka}



\begin{quote}
Before the new one: say the Telugu for an engineer, then the Telugu for a piece of chalk.
\end{quote}

\subsection*{You'll want to know: \te{పలక}}
\textbf{\te{పలక}} — \emph{palaka} — \textquotedblleft{}a slate\textquotedblright{}.

\begin{grammarlens}[title={What you've built}]
One more word for this chapter.
\end{grammarlens}

\subsection*{Your turn}
\begin{itemize}
\raggedright
  \item \emph{Say it:} \emph{palaka}
  \item \emph{Say it:} \emph{palaka}, once more
  \item \emph{Say it:} say \emph{palaka}
  \item \emph{Recall:} say the Telugu for an engineer, then the Telugu for a piece of chalk, then say \emph{palaka} again
\end{itemize}

\begin{tcolorbox}[breakable,colback=teal!4,colframe=teal!35!black,title={Before you move on}]
What does \te{పలక} mean? (\textquotedblleft{}A slate\textquotedblright{}.) Say it once more.
\end{tcolorbox}
\section[\texorpdfstring{\te{బలపం} (balapaṁ)}{బలపం (balapaṁ)}]{\te{బలపం} (balapaṁ) — a slate pencil}

\label{lesson:TE-C313-balapam}



\begin{quote}
Before the new one: say the Telugu for a piece of chalk, then the Telugu for a slate.
\end{quote}

\subsection*{You'll want to know: \te{బలపం}}
\textbf{\te{బలపం}} — \emph{balapaṁ} — \textquotedblleft{}a slate pencil\textquotedblright{}.

\begin{grammarlens}[title={What you've built}]
One more word for this chapter.
\end{grammarlens}

\subsection*{Your turn}
\begin{itemize}
\raggedright
  \item \emph{Say it:} \emph{balapaṁ}
  \item \emph{Say it:} \emph{balapaṁ}, once more
  \item \emph{Say it:} say \emph{balapaṁ}
  \item \emph{Recall:} say the Telugu for a piece of chalk, then the Telugu for a slate, then say \emph{balapaṁ} again
\end{itemize}

\begin{tcolorbox}[breakable,colback=teal!4,colframe=teal!35!black,title={Before you move on}]
What does \te{బలపం} mean? (\textquotedblleft{}A slate pencil\textquotedblright{}.) Say it once more.
\end{tcolorbox}
\section[\texorpdfstring{\te{సిరా} (sirā)}{సిరా (sirā)}]{\te{సిరా} (sirā) — ink}

\label{lesson:TE-C313-sira}



\begin{quote}
Before the new one: say the Telugu for a slate, then the Telugu for a slate pencil.
\end{quote}

\subsection*{You'll want to know: \te{సిరా}}
\textbf{\te{సిరా}} — \emph{sirā} — \textquotedblleft{}ink\textquotedblright{}.

\begin{grammarlens}[title={What you've built}]
One more word for this chapter.
\end{grammarlens}

\subsection*{Your turn}
\begin{itemize}
\raggedright
  \item \emph{Say it:} \emph{sirā}
  \item \emph{Say it:} \emph{sirā}, once more
  \item \emph{Say it:} say \emph{sirā}
  \item \emph{Recall:} say the Telugu for a slate, then the Telugu for a slate pencil, then say \emph{sirā} again
\end{itemize}

\begin{tcolorbox}[breakable,colback=teal!4,colframe=teal!35!black,title={Before you move on}]
What does \te{సిరా} mean? (\textquotedblleft{}Ink\textquotedblright{}.) Say it once more.
\end{tcolorbox}
\section[\texorpdfstring{\te{హాజరు} (hājaru)}{హాజరు (hājaru)}]{\te{హాజరు} (hājaru) — attendance}

\label{lesson:TE-C313-hajaru}



\begin{quote}
Before the new one: say the Telugu for a slate pencil, then the Telugu for ink.
\end{quote}

\subsection*{You'll want to know: \te{హాజరు}}
\textbf{\te{హాజరు}} — \emph{hājaru} — \textquotedblleft{}attendance\textquotedblright{}.

\begin{grammarlens}[title={What you've built}]
One more word for this chapter.
\end{grammarlens}

\subsection*{Your turn}
\begin{itemize}
\raggedright
  \item \emph{Say it:} \emph{hājaru}
  \item \emph{Say it:} \emph{hājaru}, once more
  \item \emph{Say it:} say \emph{hājaru}
  \item \emph{Recall:} say the Telugu for a slate pencil, then the Telugu for ink, then say \emph{hājaru} again
\end{itemize}

\begin{tcolorbox}[breakable,colback=teal!4,colframe=teal!35!black,title={Before you move on}]
What does \te{హాజరు} mean? (\textquotedblleft{}Attendance\textquotedblright{}.) Say it once more.
\end{tcolorbox}
